% =====================================================================
% 別冊 第1章 ブランチとPR編 — チームで安全に開発する
% 元原稿: 03_ブランチ・PR編.md（前半：ブランチ〜PR〜実践フロー）
% =====================================================================
\chapter{ブランチとPR編 —チームで安全に開発する—}

本書は本編『はじめてのバージョン管理』の続きです。本編では、ローカルとGitHubをつなぎ、\textbf{1本の流れ（\cmd{main}）}をアップロード／ダウンロードしました。本書では、\textbf{ブランチ}で作業を枝分かれさせ、\textbf{プルリクエスト（PR）}でレビューを挟み、\textbf{コンフリクト（競合）}を解決します。つまり「複数人・複数タスクを安全に並行させる」方法を扱います。

\begin{hintbox}
\textbf{この本の読み方}
\begin{itemize}[leftmargin=1.4em, nosep]
  \item 1人で使う分には本編で十分です。チームで開発する、または作業を安全に分けたくなったらここを読みます
  \item チームの運用ルールやPRの書き方に迷ったら → 付録A「おすすめのブランチ運用ルールとPRの書き方」へ
  \item コマンドをド忘れした人 → 『逆引きチートシート』へ
  \item エラーが出て困った人 → 『トラブルシューティング』へ
\end{itemize}
\end{hintbox}

% ---------------------------------------------------------------------
\section{なぜブランチが必要か}

これまでは \cmd{main} という1本の流れだけで作業してきました。\cmd{main} に直接編集を重ねると、こんな問題が起きます。

\begin{itemize}[leftmargin=1.4em]
  \item \textbf{未完成の変更が本流に混ざる}：書きかけのコードや原稿が \cmd{main} を壊す
  \item \textbf{複数人が同じ流れを触って衝突する}：2人が同時に \cmd{main} を編集すると \cmd{push} がぶつかる
  \item \textbf{お試しがやりにくい}：「ダメなら元に戻す」を安全に試せない
\end{itemize}

\textbf{ブランチ}を使うと、\cmd{main} から枝分かれした作業用の流れを作れます。作業はその枝の上で行い、完成してから \cmd{main} に合流（マージ）させます。こうすれば \cmd{main} は常に「動く・完成した状態」に保たれます（図\ref{fig:branch-merge}）。

\begin{numfig}{ブランチで枝分かれし、完成してから main に合流させる}{fig:branch-merge}
\begin{tikzpicture}
  \node[flownode, fill=gitGreen] (M0) {main};
  \node[flownode, fill=gitGreen, right=2.4cm of M0] (M1) {main};
  \node[flownode, fill=gitGreen, right=4.8cm of M1] (M2) {main\\（マージ後）};
  \node[flownode, fill=gitOrange, above=1.4cm of M1] (F1) {feature\\作業中};
  \node[flownode, fill=gitOrange, right=2.4cm of F1] (F2) {feature\\完成};
  \draw[flowarrow] (M0) -- (M1);
  \draw[flowarrow] (M1) -- (M2);
  \draw[flowarrow] (F1) -- (F2);
  \draw[flowarrow, dashed] (M0) to[out=70, in=200]
        node[edgelabel, midway, left] {枝分かれ} (F1);
  \draw[flowarrow, dashed] (F2) to[out=-20, in=110]
        node[edgelabel, midway, right] {マージ} (M2);
\end{tikzpicture}
\end{numfig}

% ---------------------------------------------------------------------
\section{ブランチとは}

\textbf{ブランチ（branch＝枝）}は、その名の通り履歴の枝分かれです。ある時点の \cmd{main} からコピーを取り、そこから独立して作業を進める「もう1本の作業ライン」です。

\begin{itemize}[leftmargin=1.4em]
  \item \cmd{main}（メインブランチ）… 完成した状態を保つ本流。通常ここは直接編集しない
  \item \cmd{feature/...} などの作業ブランチ … 1つのタスク専用の枝。ここで自由に編集・コミットする
\end{itemize}

各ブランチは独立しているので、作業ブランチで何をしても \cmd{main} には影響しません。\cmd{git switch} でブランチを切り替えると、作業ツリーの中身もそのブランチの内容に入れ替わります。

分け方の基本は\textbf{「1タスク＝1ブランチ」}です。「ログイン機能の追加」「READMEの修正」のように、1つの目的ごとにブランチを分けると、後のPRやレビューがきれいにまとまります。

\subsection{たとえ：倉庫で考えるブランチ}

本編では、リポジトリを書類を保管する「倉庫」にたとえました。ブランチは、この倉庫の中に置かれた\textbf{棚}だと考えると分かりやすくなります。

\begin{center}
\begin{tabularx}{\linewidth}{@{}l X@{}}
\toprule
Git & 倉庫のたとえ \\
\midrule
\cmd{main} & 全員が共有する正式な棚（本棚）。常に「完成品」だけを並べる \\
作業ブランチ & 本棚の中身を丸ごと写し取った、自分専用の作業棚 \\
そのブランチでの \cmd{commit} & 自分の作業棚に、書き直した書類の写しを並べていく \\
PR → マージ & 完成した書類を本棚に納める申請と反映 \\
\bottomrule
\end{tabularx}
\end{center}

作業ブランチを作ると、\cmd{main} の棚に並んでいる書類束（コミット）をそのまま写し取った、\textbf{自分専用の棚がもう1つ}用意されます。倉庫（リポジトリ）は1つのままで、\textbf{その中に棚が2つ並ぶ}というのが、ブランチを作るということです（図\ref{fig:branch-create}）。

\figimage[0.95]{03_branch・PR/create_branch.png}{ブランチを作ると、倉庫の中に \texttt{main} の棚と作業ブランチの棚が並ぶ}{fig:branch-create}

作業棚は本棚の写しから始まる独立した棚なので、そこで何をしても本棚（\cmd{main}）は汚れません。だから安心して書きかけの書類を置いたり、失敗した写しを捨てたりできます。

\subsection{複数人のときはどうなるか}

このたとえは、複数人での作業でこそ効いてきます。メンバーはそれぞれ、本棚から自分専用の作業棚を写し取って並行して作業します。互いの作業棚は独立しているので、他の人の手元を邪魔しません（図\ref{fig:multi-branch}）。

\begin{numfig}{各メンバーが本棚から作業棚を写し取り、完成したら本棚へ納める}{fig:multi-branch}
\begin{tikzpicture}
  \node[flownode, fill=gitGreen, text width=2.6cm] (MAIN)
        {本棚（main）\\完成品の共有棚};
  \node[flownode, fill=gitOrange, text width=3cm, above right=1.9cm and 5cm of MAIN] (A)
        {Aさんの作業棚\\feature/login};
  \node[flownode, fill=gitPurple, text width=3cm, below right=1.9cm and 5cm of MAIN] (B)
        {Bさんの作業棚\\feature/report};
  \draw[flowarrow, dashed] (MAIN.north east) to[out=55, in=210]
        node[edgelabel, pos=0.62, above left=-1pt] {写し取る} (A.west);
  \draw[flowarrow] (A.south west) to[out=250, in=30]
        node[edgelabel, pos=0.4, below right=-1pt] {PR→マージ} (MAIN.east);
  \draw[flowarrow, dashed] (MAIN.south east) to[out=-55, in=150]
        node[edgelabel, pos=0.62, below left=-1pt] {写し取る} (B.west);
  \draw[flowarrow] (B.north west) to[out=110, in=-30]
        node[edgelabel, pos=0.4, above right=-1pt] {PR→マージ} (MAIN.east);
\end{tikzpicture}
\end{numfig}

完成したら、各自が自分の作業棚の書類を順番に本棚へ納めて（マージして）いきます。ここで、AさんとBさんが「同じ書類の同じ箇所」を別々に書き換えていた場合、後から納める人のときに「どちらの内容を正とするか決められない」状態になります。これが次章の\textbf{コンフリクト（競合）}です。

逆に言えば、\textbf{触るファイルや箇所をタスクごとに分けておく}と、本棚に納める順番が前後してもぶつかりません。これがコンフリクト予防の基本です。

% ---------------------------------------------------------------------
\section{ブランチの基本操作}

作業ブランチを作って切り替えるところから始めます。

\begin{codebox}[bash]
\begin{code}
# 現在のブランチ一覧を見る（* が今いるブランチ）
git branch

# ブランチを作って、同時にそこへ切り替える
git switch -c feature/add-login

# （参考）古い書き方。同じ意味
git checkout -b feature/add-login
\end{code}
\end{codebox}

\cmd{git switch -c} の \cmd{-c}（create）は「新しく作って切り替える」という意味です。すでにあるブランチへ移動するだけなら \cmd{-c} は不要です。

\begin{codebox}[bash]
\begin{code}
# 既存のブランチへ切り替える
git switch main
git switch feature/add-login

# 不要になったブランチを削除する（マージ済みのものだけ安全に消える）
git branch -d feature/add-login
\end{code}
\end{codebox}

作業ブランチに切り替えたら、あとは本編の黄金フローと同じです。「編集 → \cmd{git add} → \cmd{git commit}」を、このブランチの上で繰り返します。

\begin{hintbox}
\textbf{ブランチ名の付け方}：厳密な決まりはありませんが、\cmd{種類/内容} の形が広く使われます。例：新機能なら \cmd{feature/add-login}、バグ修正なら \cmd{fix/typo-in-readme}。目的がひと目で分かる名前にします。おすすめの命名規則やチームの運用ルールは、付録A「おすすめのブランチ運用ルールとPRの書き方」にまとめています。
\end{hintbox}

\begin{warnbox}
作業を始める前に、\textbf{\cmd{main} を最新にしてから枝分かれ}するのが安全です（\cmd{git switch main} → \cmd{git pull} → \cmd{git switch -c 新ブランチ}）。古い \cmd{main} から分岐すると、後でコンフリクトが起きやすくなります。
\end{warnbox}

% ---------------------------------------------------------------------
\section{ブランチをリモートへ上げる}

作った作業ブランチは、ローカルにあるだけではチームに共有されません。本編の \cmd{push} でリモートにも同じブランチを作ります。

\begin{codebox}[bash]
\begin{code}
# 初回：リモートに同名のブランチを作り、宛先を記憶させる
git push -u origin feature/add-login

# 2回目以降はこれだけでOK
git push
\end{code}
\end{codebox}

\cmd{-u}（\cmd{-{}-set-upstream}）は本編と同じで、「このローカルブランチの写し先はリモートのこのブランチだ」と結び付けるオプションです。一度付ければ、以降は \cmd{git push} / \cmd{git pull} だけで今いるブランチが同期されます（図\ref{fig:push-branch}）。

\begin{numfig}{作業ブランチをリモートにも作って同期する}{fig:push-branch}
\begin{tikzpicture}
  \node[flownode, fill=gitOrange, text width=2.8cm] (LB)
        {ローカル\\feature/add-login};
  \node[flownode, fill=gitBlue, text width=2.8cm, right=8cm of LB] (RB)
        {GitHub\\origin/feature/add-login};
  \draw[flowarrow] (LB) -- (RB)
        node[edgelabel, midway, above] {git push -u origin feature/add-login};
\end{tikzpicture}
\end{numfig}

\begin{hintbox}
リモートに上げたブランチは、GitHubのリポジトリ画面のブランチ切り替えメニューから確認できます。これで、他のメンバーもあなたのブランチを見られます。
\end{hintbox}

% ---------------------------------------------------------------------
\section{プルリクエスト（PR）を作る}

\textbf{プルリクエスト（Pull Request、PR）}は、「自分の作業ブランチの変更を \cmd{main} に取り込んでください」という依頼です。ただの合流ではなく、\textbf{レビュー（変更内容の確認）を挟む入口}になります（図\ref{fig:pr-flow}）。

\begin{numfig}{PRは push からマージまでにレビューを挟む}{fig:pr-flow}
\begin{tikzpicture}
  \node[flownode, fill=gitOrange, text width=2.4cm] (B) {feature を push};
  \node[flownode, fill=gitPurple, right=2.2cm of B] (P) {PRを作成};
  \node[flownode, fill=gitBlue, right=2.2cm of P, text width=2.2cm] (R) {レビュー\\（差分を確認）};
  \node[flownode, fill=gitGreen, right=2.2cm of R] (M) {main へマージ};
  \draw[flowarrow] (B) -- (P);
  \draw[flowarrow] (P) -- (R);
  \draw[flowarrow] (R) -- (M);
\end{tikzpicture}
\end{numfig}

\noindent\textbf{PRの作成手順}

\begin{enumerate}[leftmargin=1.6em]
  \item ブランチを \cmd{push} した後、GitHubのリポジトリ画面を開きます。上部に \textbf{「Compare \& pull request」} ボタンが出るのでクリックします（出ない場合は \textbf{Pull requests} タブ → \textbf{New pull request}）。
  \item \textbf{base}（取り込み先＝通常 \cmd{main}）と \textbf{compare}（自分の作業ブランチ）が正しいか確認します。
  \item \textbf{タイトル}と\textbf{説明}を書きます。「何を・なぜ変えたか」を簡潔に書くとレビューが早く進みます。
  \item \textbf{Create pull request} をクリックします。
\end{enumerate}

\screenshot[0.62]{03_branch・PR/PR_Create_Window.png}{PR作成画面（Compare \& pull request）}
\screenshot[0.62]{03_branch・PR/Create_PR_Setting.png}{base と compare を確認し、タイトル・説明を記入して作成する}

\begin{hintbox}
\textbf{説明の書き方が伝わりやすさを決めます。} 差分を見れば「どう変えたか」は分かるので、説明には\textbf{「なぜ変えたか」}を書くのがコツです。そのまま使えるPRテンプレートと記入例を、付録A「おすすめのブランチ運用ルールとPRの書き方」に用意しています。
\end{hintbox}

PRを作ると、\textbf{以降その作業ブランチに \cmd{push} するたびに、変更が自動でPRに追加されます}。レビューの指摘を直したら、同じブランチで \cmd{add} → \cmd{commit} → \cmd{push} すればOKです。PRを作り直す必要はありません。

% ---------------------------------------------------------------------
\section{レビューとマージ}

PRの画面では、\textbf{Files changed} タブで変更の差分（追加＝緑、削除＝赤）を確認できます。チーム開発では、ここで他のメンバーがコメントや承認（Approve）を行います。

差分を見て「この変更で問題ないか」を確認します。修正が必要なら作業ブランチに戻って直し、\cmd{push} します。レビューが済んだら、PR画面の \textbf{「Merge pull request」} ボタンで \cmd{main} に取り込みます。

\screenshot[0.62]{03_branch・PR/PR_Change_Program.png}{変更内容の確認とチェック}
\screenshot[0.62]{03_branch・PR/Review_Merge.png}{レビューとマージ}

マージが済んだら、後片付けをします。

\begin{codebox}[bash]
\begin{code}
# 1. ローカルの main に戻る
git switch main

# 2. マージされた最新の main を取り込む
git pull

# 3. 役目を終えた作業ブランチを削除する（ローカル）
git branch -d feature/add-login

# 4. リモートの作業ブランチも削除する（GitHubのPR画面のボタンでも可）
git push origin --delete feature/add-login
\end{code}
\end{codebox}

マージ後にブランチを消すのは、散らかり防止のためです。GitHubのPR画面にも、マージ後に \textbf{「Delete branch」} ボタンが出ます。

次のタスクを始めるときは、また \cmd{main} を最新にしてから新しいブランチを切ります。この\textbf{「最新の \cmd{main} から枝を出す」}を守るのが、コンフリクトを減らす一番のコツです。

% ---------------------------------------------------------------------
\section{実践フロー：ブランチ運用の一周}

ここまでを1つの流れにまとめます。\textbf{この一周を繰り返す}のが、チーム開発の基本リズムです（図\ref{fig:branch-cycle}）。コンフリクトが起きたときの直し方は次章で扱います。

\begin{numfig}{ブランチ運用の一周}{fig:branch-cycle}
\begin{tikzpicture}[node distance=0.75cm]
  \node[flownode, fill=gitGreen, text width=7cm] (A) {git switch main → git pull（最新化）};
  \node[flownode, fill=gitOrange, text width=7cm, below=of A] (B) {git switch -c feature/xxx（枝を出す）};
  \node[flownode, fill=gitRed, text width=7cm, below=of B] (C) {編集 → git add → git commit};
  \node[flownode, fill=gitBlue, text width=7cm, below=of C] (D) {git push -u origin feature/xxx};
  \node[flownode, fill=gitPurple, text width=7cm, below=of D] (E) {GitHubでPR作成};
  \node[flownode, fill=gitBlue, text width=7cm, below=of E] (F) {レビュー → マージ};
  \node[flownode, fill=gitGreen, text width=7cm, below=of F] (G) {main を pull、済んだブランチを削除};
  \draw[flowarrow] (A) -- (B);
  \draw[flowarrow] (B) -- (C);
  \draw[flowarrow] (C) -- (D);
  \draw[flowarrow] (D) -- (E);
  \draw[flowarrow] (E) -- (F);
  \draw[flowarrow] (F) -- (G);
  \draw[flowarrow] (G.east) to[out=0, in=0]
        node[edgelabel, midway, right, xshift=10pt] {繰り返す} (A.east);
\end{tikzpicture}
\end{numfig}

\begin{codebox}[bash]
\begin{code}
# 1. main を最新にする
git switch main
git pull

# 2. 作業ブランチを切る
git switch -c feature/add-login

# 3. 編集して記録する（本編の黄金フロー）
git add .
git commit -m "ログイン機能を追加"

# 4. リモートへ上げる
git push -u origin feature/add-login

# 5. GitHubでPRを作成 → レビュー → マージ（ブラウザで操作）

# 6. 後片付け
git switch main
git pull
git branch -d feature/add-login
\end{code}
\end{codebox}

% ---------------------------------------------------------------------
\section*{この章のまとめ}
\addcontentsline{toc}{section}{この章のまとめ}

\begin{itemize}[leftmargin=1.4em]
  \item \textbf{\cmd{main} は直接触らず}、タスクごとに\textbf{作業ブランチ}を切って編集する
  \item 作業ブランチは \cmd{git push -u origin ブランチ名} でリモートにも共有する
  \item \textbf{PR}は「\cmd{main} に取り込んでほしい」という依頼であり、\textbf{レビューの入口}
  \item マージ後は \textbf{\cmd{main} を \cmd{pull} → 済んだブランチを削除}して片付ける
  \item \textbf{最新の \cmd{main} から枝を出す}ことがコンフリクト予防の要
\end{itemize}

次章では、同じ箇所を別々に変更したときに起きる\textbf{コンフリクト（競合）}の直し方を扱います。
